\newpage
\section{Phase-Locked Loops (PLL)}\label{sec:pll}

Los \textbf{phase-locked loops (PLL)} son circuitos de lazo cerrado que ajustan 
la fase (y por ende la frecuencia) de un oscilador controlado (\emph{VCO}/\emph{NCO}) 
para que siga a la de una señal de referencia.  
En esencia, el PLL compara la señal de entrada con una réplica generada internamente; 
si detecta un desfase, corrige al oscilador hasta que ambas coincidan.

Históricamente, los PLL se consolidaron en telecomunicaciones como bloques esenciales para la
\emph{recuperación de portadoras}, \emph{demodulación de FM/PM}, 
\emph{sincronización de fase/frecuencia} y \emph{recuperación de símbolos} 
\cite{Haykin2001,Proakis2007,Gallager2008}.  
Más allá de este uso tradicional, también se aplican al seguimiento de señales periódicas en acústica y en otros dominios donde interese estimar frecuencias dominantes 
\cite{Han2015}.

\insertarfigura{fig/PLLBlockDiagram.png}{Diagrama en bloques de un PLL clásico de segundo orden. 
	Se distinguen el detector de fase (PD), el filtro de lazo $Z(s)$ y el oscilador controlado $K_v/s$, 
	con realimentación opcional a través de un divisor $1/N$ \cite{Collins2018}.}{fig:pll:block}{1}
\insertarfigurawide{fig/PLLBank.png}{Ejemplo de un sistema con múltiples PLLs distribuidos 
	en un arreglo de antenas de RF, donde cada canal genera su propio oscilador local (LO). 
	Aunque aquí se emplea en telecomunicaciones, ilustra el concepto general de bancos de PLLs 
	trabajando en paralelo, principio que puede aplicarse también a la detección de tonos 
	en sistemas acústicos \cite{PLLBankFig}.}{fig:pll:bank}{1}
\subsection{Estructura básica e intuición}

Como muestra la Fig.~\ref{fig:pll:block}, un PLL consta de:
\begin{itemize}
	\item \textbf{Detector de fase (PD)}: mide el desfase entre la señal externa y la interna.  
	\item \textbf{Filtro de lazo $Z(s)$}: convierte la señal de error en una señal de control suave.  
	\item \textbf{Oscilador controlado (VCO/NCO)}: ajusta su frecuencia de salida según el error.  
	\item \textbf{Divisor $1/N$}: opcional en aplicaciones de síntesis de frecuencia.
\end{itemize}

El comportamiento puede explicarse así:  
si la señal de entrada tiene una frecuencia ligeramente superior a la del oscilador, 
el detector de fase produce un error positivo que acelera al VCO.  
Si es más lenta, el error lo frena.  
El proceso se repite hasta que el VCO acompaña de forma estable a la señal externa.  
De este modo, el PLL actúa como un seguidor dinámico de fase y frecuencia.
\balance
\subsection{Modelo lineal}

Cuando el lazo ya está cerca del enganche, el análisis puede simplificarse a un modelo lineal.  
El error de fase $\phi_e = \phi_{\text{ref}} - \phi_{\text{vco}}$ obedece:
\begin{equation}
	\Phi_e(s) \;=\; \frac{1}{1 + K_d K_v \, \tfrac{Z(s)}{s}}\, \Phi_{\text{ref}}(s),
\end{equation}
donde $K_d$ es la ganancia del detector de fase, $K_v$ la sensibilidad del VCO y $Z(s)$ el filtro.  

Si el filtro es proporcional-integral (PI):
\begin{equation}
	Z(s) = K_p + \frac{K_i}{s},
\end{equation}
el lazo cerrado se comporta como un sistema de segundo orden con:
\[
\omega_n = \sqrt{K_d K_v K_i}, \qquad 
\zeta = \tfrac{1}{2}\sqrt{\tfrac{K_d K_v K_p}{K_i}},
\]
donde $\omega_n$ es la frecuencia natural del lazo y $\zeta$ el factor de amortiguamiento 
\cite{Collins2018,ADI_MT086}.  
Estos parámetros controlan la rapidez con la que el PLL se engancha, su estabilidad y si aparecen sobreoscilaciones.

\subsection{Rangos prácticos y ruido}

Tres magnitudes prácticas caracterizan el desempeño de un PLL \cite{Collins2018}:
\begin{itemize}
	\item \textbf{Lock range}: rango de frecuencias donde el lazo puede mantenerse enganchado.  
	\item \textbf{Pull-in range}: rango desde el cual el lazo es capaz de lograr enganche inicial.  
	\item \textbf{Ancho de banda $B_L$}: cuánto se deja pasar de variaciones rápidas y de ruido.
\end{itemize}
Un ancho de banda grande hace al lazo rápido, pero lo expone más al ruido; 
uno estrecho lo protege, pero lo hace lento.

\subsection{Implementación digital (DPLL)}

En sistemas digitales, el VCO se reemplaza por un \emph{NCO} (Numerically Controlled Oscillator) 
y todos los bloques se discretizan.  
El detector de fase puede implementarse, por ejemplo, mezclando la señal de entrada con una referencia en cuadratura ($I$ y $Q$) y calculando:
\[
e_\phi[n] = \operatorname{atan2}(Q[n], I[n]).
\]

El filtro PI en forma discreta acumula el error de fase:
\begin{align}
	u[n] &= u[n-1] + K_p \, e_\phi[n] + K_i \sum_{k=0}^{n} e_\phi[k], \\
	\phi_{\text{nco}}[n+1] &= \phi_{\text{nco}}[n] + K_v \, u[n].
\end{align}

Aquí:
\begin{itemize}
	\item $K_p$ define la respuesta rápida a cambios instantáneos,  
	\item $K_i$ corrige errores de fase acumulados a largo plazo,  
	\item $K_v$ traduce la señal de control en cambio de frecuencia del NCO.
\end{itemize}

En la práctica, este cálculo se implementa con sumadores, multiplicadores y acumuladores digitales, 
por lo que es especialmente amigable a hardware como FPGAs.  
Técnicas como Tustin o bilineal se usan para discretizar $K_p, K_i$ a partir del diseño analógico.

\subsection{Banco de PLLs}
Un ejemplo de este enfoque se observa en la Fig.~\ref{fig:pll:bank}. 
En lugar de un único lazo, puede emplearse un banco de PLLs en paralelo, cada uno sintonizado a una frecuencia $f_k$.  
Cada PLL entrega tres señales clave:
\begin{itemize}
	\item Un indicador de enganche (lock).  
	\item La estimación de fase/frecuencia en su canal.  
	\item La potencia coherente asociada a esa frecuencia.
\end{itemize}
Si un conjunto de PLLs (ej. los que cubren $f_0, 2f_0, 3f_0$) se enganchan al mismo tiempo, 
se infiere la presencia de una señal con ese patrón armónico.  
Esto equivale a un análisis espectral adaptativo con muy baja latencia y es factible en hardware embebido 
\cite{Han2015}.

\subsection{Ventajas y limitaciones}

\textbf{Ventajas}: bajo costo computacional, buena selectividad, fácil de implementar en FPGA/DSP.  
\textbf{Limitaciones}: rango de captura limitado, no separa dos fuentes con la misma frecuencia, 
y puede perder enganche en ambientes con multitrayectoria severa \cite{Collins2018}.

\subsection{Relevancia y aplicaciones}

El uso de \textbf{PLLs} en este trabajo no responde únicamente a una cuestión técnica, 
sino a la necesidad de \textbf{detectar y seguir componentes periódicas} dentro de señales ruidosas.  
En el análisis de audio, esto es relevante porque muchas fuentes de interés —incluyendo insectos, 
motores eléctricos o vocalizaciones humanas— se caracterizan por su periodicidad.  
El PLL permite extraer esa frecuencia dominante en tiempo real con bajo costo computacional, 
incluso cuando el entorno introduce variaciones o interferencias.

El problema que resuelven es claro: \textit{¿cómo identificar y mantener la referencia de una señal 
	cuando ésta se encuentra sumergida en ruido o compite con otras fuentes?}.  
A diferencia de un filtrado espectral general, que entrega un panorama global pero con mayor latencia, 
un PLL está optimizado para \textbf{seguir dinámicamente una frecuencia puntual} y entregar medidas 
instantáneas de fase y potencia coherente.  
Por ello resulta más adecuado en contextos donde interesa la estabilidad y rapidez de respuesta 
(monitorización en drones, sistemas de control embebido, comunicaciones en tiempo real).

En cuanto a aplicaciones más allá de los mosquitos, el mismo principio se utiliza en la 
\textbf{sincronización de inversores de potencia con la red}, en la \textbf{recuperación de portadoras} 
en receptores de radio o satélite, en \textbf{GPS para alinear relojes locales}, 
e incluso en \textbf{biomedicina} para seguir ritmos cardiacos o respiratorios.  
En contextos acústicos, podría servir para identificar patrones de maquinaria en mantenimiento predictivo 
o para seguir frecuencias de habla en sistemas de mejora auditiva.

Así, el análisis con PLL se justifica no sólo como técnica histórica de telecomunicaciones, 
sino como un bloque versátil que puede compararse con métodos modernos:  
ofrece interpretabilidad y eficiencia donde se requiere simplicidad, 
mientras que enfoques más complejos (ej. \emph{deep learning}) se destinan a escenarios con mayor variabilidad.  
Ambos enfoques no se excluyen, sino que pueden combinarse para formar sistemas híbridos.



